\documentclass[12pt,letterpaper]{exam}
\usepackage[utf8]{inputenc}
\usepackage[T1]{fontenc}
\usepackage[width=8.50in, height=11.00in, left=0.50in, right=0.50in, top=0.50in, bottom=0.50in]{geometry}

\usepackage{libertine}
\usepackage{multicol}
\usepackage[shortlabels]{enumitem}

\usepackage{booktabs}
\usepackage[table]{xcolor}

\usepackage{amssymb}
\usepackage{amsthm}
\usepackage{mathtools}
\usepackage{bbm}

\usepackage{hyperref}
\usepackage{graphicx}
%\usepackage{wrapfig}
%\usepackage{capt-of}
\usepackage{tikz}
%\usepackage{pgfplots}
%\usetikzlibrary{shapes,arrows,positioning,patterns}
%\usepackage{pythonhighlight}

\newcommand\chapter{4}
\renewcommand{\thequestion}{\textbf{\chapter.\arabic{question}}}
\renewcommand{\questionlabel}{\thequestion}

%%%%%%%%%%%%%%%%%%%%%%%%%%%%%%%%%%%%%%%%%%%%%%%%%%%%%%%%%%%%%%%%%%
\newcommand{\class}{STAT 601} % This is the name of the course 
\newcommand{\assignmentname}{Homework \# \chapter} % 
\newcommand{\authorname}{Hosley, Brandon} % 
\newcommand{\workdate}{\today} % 
\printanswers % this includes the solutions sections
%%%%%%%%%%%%%%%%%%%%%%%%%%%%%%%%%%%%%%%%%%%%%%%%%%%%%%%%%%%%%%%%%%



\begin{document}
\pagestyle{plain}
\thispagestyle{empty}
\noindent
 
%%%%%%%%%%%%%%%%%%%%%%%%%%%%%%%%%%%%%%%%%%%%%%%%%%%%%%%%%%%%%%%%%%%%%%%%%%%%%%%%%%%
\noindent
\begin{tabular*}{\textwidth}{l @{\extracolsep{\fill}} r @{\extracolsep{10pt}} l}
	\textbf{\class} & \textbf{\authorname}  &\\ %Your name here instead, obviously 
	\textbf{\assignmentname } & \textbf{\workdate} & \\
\end{tabular*}\\ 
\rule{\textwidth}{2pt}
%%%%%%%%%%%%%%%%%%%%%%%%%%%%%%%% HEADER %%%%%%%%%%%%%%%%%%%%%%%%%%%%%%%%%%%%%%%%%%%

\begin{questions}

	%4.4
	\setcounter{question}{3}
	\question 
	A pdf is defined by
	
	\[
	f(x,y)=
	\begin{cases}
		C(x + 2y) & \text{if } 0 < y < 1 \text{ and } 0 < x < 2 \\
		0 & \text{otherwise}.
	\end{cases}
	\]
	
	\begin{parts}
		\part Find the value of \(C\).
		\part Find the marginal distribution of \(X\).
		\part Find the joint cdf of \(X\) and \(Y\).
		\part Find the pdf of the random variable \(Z = 9/(X + 1)^2\).
	\end{parts}
	
	\begin{solution}
		\begin{parts}
			\part
			If \(f(x,y)\) is a valid pdf then,
			\begin{align*}
				\int f(x,y) = 1
				&= \int_{0}^{2} \int_{0}^{1} C(x+2y) \,dy\,dx \\
				&= \int_{0}^{2} \bigg[Cxy+Cy^2\bigg]_0^1 \,dx \\
				&= \int_{0}^{2} Cx+C \,dx \\
				&= \left. Cx+\frac{Cx^2}{2} \right|_{0}^{2} \\
				&= 4C, \\
			\end{align*}
			thus we conclude that \(C = \frac{1}{4}\).
			
			\part
			We obtain the marginal by integrating out all other variables,
			\begin{align*}
				f_X(x)
				=\ \int_{0}^{1} \frac{1}{4} (x+2y) \,dy
				=\ \left. \frac{1}{4}xy + \frac{1}{4}y^2 \right|_0^1
				=\ \frac{1}{4}(x+1).
			\end{align*}

			\part
			For \(0<x<2\) and \(0<y<1\),
			\begin{align*}
				F(x,y)
				&= \int_{0}^{x} \int_{0}^{y} \frac{1}{4}x + \frac{1}{2}y \,dy \,dx \\
				&= \int_{0}^{x} \frac{1}{4}xy + \frac{1}{4}y^2 \,dx \\
				&= \frac{1}{8}x^2y + \frac{1}{4}xy^2 \\
				&= \frac{1}{4}xy \left(\frac{1}{2}x + y \right). \\
			\end{align*}
			
			The unintuitive extension of the distribution beyond joint feasibility,
			
			\(\frac{1}{8}x^2 + \frac{1}{4}x\ \) for \(1 \leq y\) and \(0<x<2\), and
			
			\(\frac{1}{4}y^2 + \frac{1}{8}y\ \) for \(2 \leq x\) and \(0<y<1\).
			
			\part
			
			\(g^{-1}(x) = \sqrt{\frac{9}{z}} - 1\)
			
			\begin{align*}
				f_Z(z)
				&= \frac{1}{4} \left( \sqrt{\frac{9}{z}} \right) \left| \frac{d}{dz} g^{-1}(x) \right| \\
				&= \frac{1}{4} \left( \frac{3}{\sqrt{z}} \right) \left( \frac{3}{2} z^{-\frac{3}{2}} \right) \\
				&= \frac{9}{8} z^{-2}
			\end{align*}
			
			Using \(g^{-1}(x)\) we calculate the support to be over \(1<z<9\).
			
		\end{parts}		
	\end{solution}
	%%%%%%%%%%%%%%%%%%%%%%%%%%%%%%%%%%%%%%%%%%%%%%%%%%%%%%%%%%%%%

	%4.5
	\question 
	\begin{parts}
		\part Find \(P(X>\sqrt{Y})\) if \(X\) and \(Y\) are jointly distributed with pdf
			\[f(x,y) = x + y, \quad 0 \leq x \leq 1, \quad 0 \leq y \leq 1.\]
		\part Find \(P(X^2 < Y < X)\) if \(X\) and \(Y\) are jointly distributed with pdf
			\[f(x,y) = 2x, \quad 0 \leq x \leq 1, \quad 0 \leq y \leq 1.\]
	\end{parts}
	
	\begin{solution}
		\begin{parts}
			\part
			
			\begin{align*}
				P(X > \sqrt{Y})
				&= \int_{0}^{1} \int_{\sqrt{y}}^{1} x + y \,dx \,dy \\
				&= \int_{0}^{1} \left[ \frac{1}{2}x^2 + xy \right]_{\sqrt{y}}^{1} \,dy \\
				&= \int_{0}^{1} \left[ \left(\frac{1}{2} + y\right) - \left(\frac{1}{2}y + y^{3/2}\right) \right] \,dy \\
				&= \int_{0}^{1} \frac{1}{2} + \frac{1}{2}y - y^{3/2} \,dy \\
				&= \left. \frac{1}{2}y + \frac{1}{4}y^2 - \frac{2}{5}y^{5/2} \right|_{0}^{1} \\
				&= \frac{1}{2} + \frac{1}{4} - \frac{2}{5} \\
				&= \frac{7}{20} \\
			\end{align*}
			
			\part
			
			\begin{align*}
				P(X^2 < Y < X)
				&= \int_{0}^{1} \int_{X^2}^{X} 2x \,dy \,dx \\
				&= \int_{0}^{1} \left[2xy\right]_{X^2}^{X} \,dx \\
				&= \int_{0}^{1} 2x^2 - 2x^3 \,dx \\
				&= \left. \frac{2}{3}x^3 - \frac{1}{2}x^4 \right|_{0}^{1} \\
				&= \frac{2}{3} - \frac{1}{2} \\
				&= \frac{1}{6} \\
			\end{align*}			
			
		\end{parts}
	\end{solution}
	%%%%%%%%%%%%%%%%%%%%%%%%%%%%%%%%%%%%%%%%%%%%%%%%%%%%%%%%%%%%%

	%4.7
	\setcounter{question}{6}
	\question 
	A woman leaves for work between 8 AM and 8:30 AM and takes between 40 and 50 
	minutes to get there. Let the random variable \(X\) denote her time of departure, and
	the random variable \(Y\) the travel time. Assuming that these variables are independent
	and uniformly distributed, find the probability that the woman arrives at work before
	9 AM.
	
	\begin{solution}
		Support is over \(0<x<30\), \(40<y<50\). The marginal probabilities for each part is		
		\(f(x) = \frac{1}{30} \mathbbm{1}_{(0,30)}(x)\) 
		and
		\(f(y) = \frac{1}{10} \mathbbm{1}_{(40,50)}(y)\).
		\begin{align*}
			P(X+Y<60)
			&= P(X<60-Y) \\
			&= \int_{40}^{50} \int_{0}^{60-y} \frac{1}{30} \mathbbm{1}(x) \frac{1}{10} \mathbbm{1}(y) \,dx \,dy \\
			&= \int_{40}^{50} \left. \frac{1}{30}x \right|_{0}^{60-y} \frac{1}{10} \mathbbm{1}(y) \,dy \\
			&= \int_{40}^{50} \frac{1}{5} - \frac{y}{300} \,dy \\
			&= \left. \frac{y}{5} - \frac{y^2}{300} \right|_{40}^{50} \\
			&= \left(10 - \frac{25}{6}\right) - \left(8 - \frac{16}{6}\right) \\
			&= 2 - \frac{9}{2} \\
			&= \frac{1}{2} .
		\end{align*}
		
	\end{solution}
	%%%%%%%%%%%%%%%%%%%%%%%%%%%%%%%%%%%%%%%%%%%%%%%%%%%%%%%%%%%%%
	\clearpage
	
	%4.9
	\setcounter{question}{8}
	\question 
	Prove that if the joint cdf of \(X\) and \(Y\) satisfies
		\[F_{X,Y} (x,y) = F_X(x)F_Y(y),\]
	then for any pair of intervals \((a, b)\), and \((c, d)\),
		\[P(a \leq X \leq b, c \leq Y \leq d) = P(a \leq X \leq b)P(c \leq Y \leq d).\]
	
	\begin{solution}
		
	\end{solution}
	%%%%%%%%%%%%%%%%%%%%%%%%%%%%%%%%%%%%%%%%%%%%%%%%%%%%%%%%%%%%%

	%4.10
	\question 
	The random pair \((X,Y)\) has the distribution
	
	\begin{center}
		\def\arraystretch{1.25}
		\begin{tabular}{cc|ccc}
			& & & \(X\) & \\
			& & 1 & 2 & 3 \\
			\hline
			& 2 & \(\frac{1}{12}\) & \(\frac{1}{6}\) & \(\frac{1}{12}\) \\
			\(Y\) & 3 & \(\frac{1}{6}\)  & \(0\) & \(\frac{1}{6}\) \\
			& 4 & \(0\) & \(\frac{1}{3}\) &  \(0\)
		\end{tabular}
	\end{center}
	
	
	\begin{parts}
		\part Show that \(X\) and \(Y\) are dependent.
		\part Give a probability table for random variables \(U\) and \(V\) 
			that have the same marginals as \(X\) and \(Y\) but are independent.
	\end{parts}
	
	\begin{solution}
		\begin{parts}
			\part
			If the variables are independent then \(P(Y|X) = P(Y)\).
			
			However, \(P(Y=4|X=1) = 0 \neq P(Y=4) = \frac{1}{4}\).
			
			\part \phantom{text}
			
			\begin{center}
				\def\arraystretch{1.25}
				\begin{tabular}{cc|ccc}
					& & & \(U\) & \\
					& & 1 & 2 & 3 \\
					\hline
					& 2 & \(\frac{1}{12}\) & \(\frac{1}{6}\) & \(\frac{1}{12}\) \\
					\(V\) & 3 & \(\frac{1}{12}\)  & \(\frac{1}{6}\) & \(\frac{1}{12}\) \\
					& 4 &\(\frac{1}{12}\) & \(\frac{1}{6}\) &  \(\frac{1}{12}\)
				\end{tabular}
			\end{center}
						
		\end{parts}
	\end{solution}
	%%%%%%%%%%%%%%%%%%%%%%%%%%%%%%%%%%%%%%%%%%%%%%%%%%%%%%%%%%%%%

	%4.12
	\setcounter{question}{11}
	\question 
	If a stick is broken at random into three pieces, what is the probability that the pieces
	can be put together in a triangle? (See Gardner 1961 for a complete discussion of this
	problem.)
	
	\begin{solution}
		
	\end{solution}
	%%%%%%%%%%%%%%%%%%%%%%%%%%%%%%%%%%%%%%%%%%%%%%%%%%%%%%%%%%%%%

	%4.19
	\setcounter{question}{18}
	\question 
	
	\begin{parts}
		\part Let \(X_1\) and \(X_2\) be independent n\((0, 1)\) random variables. 
			Find the pdf of \((X_1 - X_2)^2/2\).
		\part If \(X_i, i = 1, 2,\) are independent gamma\((\alpha_i, 1)\) random variables, 
			find the marginal distributions of \(X_1/(X_1 + X_2)\) and \(X_2/(X_1 + X_2)\).
	\end{parts}
	
	\begin{solution}
		
		\begin{tikzpicture}
			
			% Draw rectangle segments and label them
			\draw[fill=red,  opacity=0.2] (0,0) rectangle (2,0.5) node[midway] {A};
			\draw[fill=blue, opacity=0.2] (2,0) rectangle (4,0.5) node[midway] {B};
			\draw[fill=green,  opacity=0.2] (4,0) rectangle (6,0.5) node[midway] {C};
			\draw[fill=yellow,  opacity=0.2] (6,0) rectangle (8,0.5) node[midway] {D};
			
		\end{tikzpicture}
		
	\end{solution}
	%%%%%%%%%%%%%%%%%%%%%%%%%%%%%%%%%%%%%%%%%%%%%%%%%%%%%%%%%%%%%

	%4.23
	\setcounter{question}{22}
	\question 
	For \(X\) and \(Y\) as in Example 4.3.3, find the distribution of \(XY\) by making the transformations given in (a) and (b) and integrating out \(V\).
	
	\begin{parts}
		\part \(U = XY ,\ V = Y\)
		\part \(U = XY ,\ V = X/Y\)
	\end{parts}
	
	\begin{solution}
		
	\end{solution}
	%%%%%%%%%%%%%%%%%%%%%%%%%%%%%%%%%%%%%%%%%%%%%%%%%%%%%%%%%%%%%

	%4.24
	\question
	Let \(X\) and \(Y\) be independent random variables with 
	\(X\sim\)gamma\((r,1)\) and \(Y\sim\)gamma\((s,1)\). 
	Show that \(Z_1 = X + Y\) and \(Z_2 = X/(X + Y )\) are independent, 
	and find the distribution of each. (\(Z_1\) is gamma and \(Z_2\) is beta.) 
	
	\begin{solution}
		
	\end{solution}
	%%%%%%%%%%%%%%%%%%%%%%%%%%%%%%%%%%%%%%%%%%%%%%%%%%%%%%%%%%%%%

	%4.30
	\setcounter{question}{30}
	\question 
	Suppose the distribution of \(Y\), conditional on \(X = x\), is n\((x, x^2)\) 
	and that the marginal distribution of \(X\) is uniform\((0, 1)\).
	
	\begin{parts}
		\part Find \(E\,Y\) , Var \(Y\) , and Cov\((X,Y)\).
		\part Prove that \(Y/X\) and \(X\) are independent.
	\end{parts}
	
	\begin{solution}
		
	\end{solution}
	%%%%%%%%%%%%%%%%%%%%%%%%%%%%%%%%%%%%%%%%%%%%%%%%%%%%%%%%%%%%%

	%4.31
	\question 
	Suppose that the random variable \(Y\) has a binomial distribution with \(n\) trials and success probability \(X\), where \(n\) is a given constant and \(X\) is a uniform\((0, 1)\) random variable.
	
	\begin{parts}
		\part Find \(E\,Y\) and Var \(Y\).
		\part Find the joint distribution of \(X\) and \(Y\).
		\part Find the marginal distribution of \(Y\).
	\end{parts}
	
	\begin{solution}
		
	\end{solution}
	%%%%%%%%%%%%%%%%%%%%%%%%%%%%%%%%%%%%%%%%%%%%%%%%%%%%%%%%%%%%%

	%4.41
	\setcounter{question}{40}
	\question 
	Show that any random variable is uncorrelated with a constant.
	
	\begin{solution}
		
	\end{solution}
	%%%%%%%%%%%%%%%%%%%%%%%%%%%%%%%%%%%%%%%%%%%%%%%%%%%%%%%%%%%%%

	%4.45
	\setcounter{question}{44}
	\question 
	Show that if \((X,Y)\sim\) bivariate normal
	\((\mu_X, \mu_Y , \sigma_X^2 , \sigma_Y^2 , \rho)\), then the following are true.
	
	\begin{parts}
		\part The marginal distribution of \(X\) is n\((\mu_X, \sigma_X^2 )\) a
		nd the marginal distribution of \(Y\) is n\((\mu_Y , \sigma_Y^2 )\).
		\part The conditional distribution of Y given \(X = x\) is
			\[n(\mu_Y + \rho(\sigma_Y /\sigma_X)(x-\mu_X), \sigma_Y^2 (1-\rho^2)).\]
		\part For any constants \(a\) and \(b\), the distribution of \(aX + bY\) is
			\[n(a\mu_X + b\mu_Y , a^2\sigma_X^2 + b^2\sigma_Y^2 + 2ab\rho\sigma_X\sigma_Y ).\]
	\end{parts}
	
	\begin{solution}
		
	\end{solution}
	%%%%%%%%%%%%%%%%%%%%%%%%%%%%%%%%%%%%%%%%%%%%%%%%%%%%%%%%%%%%%

	%4.52	
	\setcounter{question}{51}
	\question 
	Bullets are fired at the origin of an \((x, y)\) coordinate system, and the point hit, say \((X, Y )\), is a random variable. The variables \(X\) and \(Y\) are taken to be independent n\((0, 1)\) random variables. If two bullets are fired independently, what is the distribution of the distance between them?
	
	\begin{solution}
		
	\end{solution}
	%%%%%%%%%%%%%%%%%%%%%%%%%%%%%%%%%%%%%%%%%%%%%%%%%%%%%%%%%%%%%

\end{questions}
\end{document}
